\documentclass{article}
\usepackage[UTF8]{ctex}
\usepackage{amsmath,amssymb,amsfonts}
\usepackage{enumitem}
\usepackage{geometry}
\geometry{a4paper, margin=1in}
% === 错题本排版优化（不参与题目标识） ===
\setlength{\parindent}{0pt}
\setlength{\parskip}{0.35em plus 0.10em minus 0.08em}
\setlength{\emergencystretch}{3em}
\setlist{itemsep=0.25em, topsep=0.40em, parsep=0pt}
\allowdisplaybreaks[2]
% === 错题本排版优化结束 ===
\begin{document}
\section*{哈尔滨工程大学2024年高等代数真题}
\section*{77．哈尔滨工程大学2024年高等代数试卷}
一．填空题．\\

\subsection{题目1}
1．设 \(a \neq b, n\) 阶行列式 \(\left|\begin{array}{cccccc}a+b & a b & 0 & \cdots & 0 & 0 \\ 1 & a+b & a b & \cdots & 0 & 0 \\ 0 & 1 & a+b & \cdots & 0 & 0 \\ \vdots & \vdots & \vdots & & \vdots & \vdots \\ 0 & 0 & 0 & \cdots & a+b & a b \\ 0 & 0 & 0 & \cdots & 1 & a+b\end{array}\right|\) 的值为 \(\_\_\_\_\)．\\
\(\_\_\_\_\)

\subsection{题目2}
2．若 3 阶可逆阵 \(A\) 交换 1,2 行得矩阵 \(B\) ，再将矩阵 \(B\) 的第 2 行加到第 3 行上得到矩阵 \(C\) ，则满足\(P A=C\) 的矩阵 \(P=\) \(\_\_\_\_\)．

\subsection{题目3}
3．若二次型 \(f\left(x_{1}, x_{2}, x_{3}\right)=t\left(x_{1}^{2}+x_{2}^{2}+x_{3}^{2}\right)+2 x_{1} x_{2}-2 x_{2} x_{3}\) 正定，则参数 \(t\) 应满足 \(\_\_\_\_\)\\
\(\_\_\_\_\)．．\\
\(\_\_\_\_\)

\subsection{题目4}
4．设 \(A, B\) 均为 \(m \times n\) 矩阵，且 \(\operatorname{rank} A=\operatorname{rank} B=1, W_{1}\) 与 \(W_{2}\) 分别为齐次线性方程组 \(A X=0\) 和 \(B X=0\) 的解空间，且 \(W_{1} \neq W_{2}\) ，则 \(\operatorname{dim}\left(W_{1} \cap W_{2}\right)=\) \(\_\_\_\_\)．

\subsection{题目5}
5．设 3 阶对称阵 \(A\) 满足 \(A^{2}+2 A-3 E=O\) ，其中 \(E\) 为单位阵，且 \(A\) 与 \(B=\left(\begin{array}{lll}2 & 3 & 3 \\ 3 & 2 & 3 \\ 3 & 3 & 2\end{array}\right)\) 合同，则 \(|A+2 E|=\) \(\_\_\_\_\)．
\end{document}
